%!TEX root = ../main.tex
\chapter{Monitoring des Mod\`eles en Production}
\label{ch:monitoring}

\section{Introduction}

D\'eployer un mod\`ele de Machine Learning en production n'est pas la fin du processus~: c'est le d\'ebut d'une nouvelle phase critique. Les mod\`eles se d\'egradent in\'evitablement avec le temps, car les donn\'ees du monde r\'eel \'evoluent. Le monitoring permet de d\'etecter ces d\'egradations avant qu'elles n'impactent les utilisateurs.

\begin{attention}[title=\textbf{Pourquoi les mod\`eles se d\'egradent-ils~?}]
Un mod\`ele entra\^in\'e capture une relation $f: X \to Y$ valide au moment de l'entra\^inement. Cette relation peut devenir obsol\`ete pour plusieurs raisons~:
\begin{itemize}
    \item \textbf{Les donn\'ees changent}~: le profil des utilisateurs \'evolue, de nouveaux comportements apparaissent.
    \item \textbf{Le contexte change}~: saisonnalit\'e, \'ev\'enements \'economiques, changements r\'eglementaires.
    \item \textbf{Le syst\`eme en amont change}~: modification d'un capteur, changement de format de donn\'ees.
    \item \textbf{Les labels changent}~: ce qui \'etait consid\'er\'e comme positif hier ne l'est plus aujourd'hui.
\end{itemize}
Sans monitoring, ces d\'egradations passent inaper\c{c}ues pendant des semaines ou des mois.
\end{attention}

%% ----------------------------------------------------------------
\section{Types de drift}
\label{sec:drift-types}

Le \emph{drift} (d\'erive) d\'esigne tout changement dans la distribution des donn\'ees ou dans la relation entre les entr\'ees et les sorties d'un mod\`ele.

\begin{definition}[Covariate Shift]
Le \emph{covariate shift} (d\'erive des covariables) survient lorsque la distribution des variables d'entr\'ee $P(X)$ change, tandis que la relation conditionnelle $P(Y|X)$ reste identique. Par exemple, un mod\`ele de cr\'edit entra\^in\'e sur des clients de 25--55 ans qui re\c{c}oit soudainement des demandes de clients de 18--25 ans.
\[
P_{\text{train}}(X) \neq P_{\text{prod}}(X), \quad P_{\text{train}}(Y|X) = P_{\text{prod}}(Y|X)
\]
\end{definition}

\begin{definition}[Prior Probability Shift]
Le \emph{prior probability shift} (d\'erive de la distribution cible) survient lorsque la distribution de la variable cible $P(Y)$ change. Par exemple, un classificateur de spam entra\^in\'e avec 10\,\% de spam qui fait face \`a 40\,\% de spam en production.
\[
P_{\text{train}}(Y) \neq P_{\text{prod}}(Y), \quad P_{\text{train}}(X|Y) = P_{\text{prod}}(X|Y)
\]
\end{definition}

\begin{definition}[Concept Drift]
Le \emph{concept drift} (d\'erive de concept) est le cas le plus s\'ev\`ere~: la relation entre les entr\'ees et la cible change fondamentalement.
\[
P_{\text{train}}(Y|X) \neq P_{\text{prod}}(Y|X)
\]
Exemple~: un mod\`ele de recommandation dont le comportement des utilisateurs change radicalement apr\`es une crise \'economique.
\end{definition}

\begin{figure}[htbp]
\centering
\begin{tikzpicture}[
    box/.style={rectangle, draw, rounded corners, minimum width=3.5cm,
        minimum height=1.5cm, align=center, font=\small},
    arrow/.style={-{Stealth[length=2.5mm]}, thick},
    title/.style={font=\bfseries\small}
]

% Covariate shift
\node[box, fill=blue!10] (cov) at (0,0) {Covariate Shift\\$P(X)$ change\\$P(Y|X)$ stable};

% Prior probability shift
\node[box, fill=green!10] (prior) at (5.5,0) {Prior Probability\\Shift\\$P(Y)$ change\\$P(X|Y)$ stable};

% Concept drift
\node[box, fill=red!10] (concept) at (11,0) {Concept Drift\\$P(Y|X)$ change};

% Severity arrow
\draw[arrow, very thick, gray!60] (0,-1.8) -- (11,-1.8)
    node[midway, below, font=\small\bfseries, text=gray!80] {S\'ev\'erit\'e croissante};

% Labels
\node[title, above=0.3cm of cov] {Mod\'er\'e};
\node[title, above=0.3cm of prior] {Significatif};
\node[title, above=0.3cm of concept] {Critique};

\end{tikzpicture}
\caption{Les trois types de drift class\'es par s\'ev\'erit\'e.}
\label{fig:drift-types}
\end{figure}

%% ----------------------------------------------------------------
\section{Formulation math\'ematique de la d\'etection de drift}
\label{sec:drift-math}

\begin{formules}[title=\textbf{M\'etriques de d\'etection de drift}]
\textbf{Divergence de Kullback-Leibler (KL)}~:
\[
D_{\mathrm{KL}}(P \| Q) = \sum_{x} P(x) \ln \frac{P(x)}{Q(x)}
\]
La divergence KL est \textbf{asym\'etrique}~: $D_{\mathrm{KL}}(P \| Q) \neq D_{\mathrm{KL}}(Q \| P)$. Elle mesure la quantit\'e d'information perdue lorsqu'on utilise $Q$ pour approximer $P$.

\medskip
\textbf{Population Stability Index (PSI)}~:
\[
\mathrm{PSI} = \sum_{i=1}^{k} (p_i - q_i) \ln \frac{p_i}{q_i}
\]
o\`u $p_i$ et $q_i$ sont les proportions dans le $i$-\`eme intervalle (\emph{bin}) des distributions de r\'ef\'erence et courante. R\`egles d'interpr\'etation~:
\begin{itemize}
    \item $\mathrm{PSI} < 0.1$~: pas de changement significatif
    \item $0.1 \leq \mathrm{PSI} < 0.2$~: changement mod\'er\'e, \`a surveiller
    \item $\mathrm{PSI} \geq 0.2$~: changement significatif, action requise
\end{itemize}

\medskip
\textbf{Test de Kolmogorov-Smirnov (KS)}~:
\[
D_n = \sup_x |F_{\text{ref}}(x) - F_{\text{cur}}(x)|
\]
o\`u $F_{\text{ref}}$ et $F_{\text{cur}}$ sont les fonctions de r\'epartition empiriques des distributions de r\'ef\'erence et courante. On rejette l'hypoth\`ese nulle (<<~pas de drift~>>) si la \emph{p-value} est inf\'erieure au seuil $\alpha$ (typiquement 0.05).
\end{formules}

\begin{codeblock}[title=\textbf{Impl\'ementation des m\'etriques de drift}]
\begin{minted}{python}
import numpy as np
from scipy import stats


def kl_divergence(p: np.ndarray, q: np.ndarray,
                  epsilon: float = 1e-10) -> float:
    """Divergence KL entre deux distributions discretes."""
    p = np.clip(p, epsilon, None)
    q = np.clip(q, epsilon, None)
    p = p / p.sum()
    q = q / q.sum()
    return float(np.sum(p * np.log(p / q)))


def psi(reference: np.ndarray, current: np.ndarray,
        n_bins: int = 10) -> float:
    """Population Stability Index."""
    breakpoints = np.linspace(
        min(reference.min(), current.min()),
        max(reference.max(), current.max()),
        n_bins + 1,
    )
    ref_counts = np.histogram(reference, bins=breakpoints)[0]
    cur_counts = np.histogram(current, bins=breakpoints)[0]

    # Eviter les zeros
    ref_pct = (ref_counts + 1) / (len(reference) + n_bins)
    cur_pct = (cur_counts + 1) / (len(current) + n_bins)

    return float(np.sum((cur_pct - ref_pct) * np.log(cur_pct / ref_pct)))


def ks_test(reference: np.ndarray, current: np.ndarray,
            alpha: float = 0.05) -> dict:
    """Test de Kolmogorov-Smirnov pour detecter le drift."""
    statistic, p_value = stats.ks_2samp(reference, current)
    return {
        "statistic": float(statistic),
        "p_value": float(p_value),
        "drift_detected": p_value < alpha,
    }
\end{minted}
\end{codeblock}

%% ----------------------------------------------------------------
\section{Evidently AI~: rapports de monitoring}
\label{sec:evidently}

Evidently est une biblioth\`eque open-source Python qui g\'en\`ere des rapports de monitoring pour les mod\`eles ML. Elle permet de d\'etecter le drift, d'analyser la qualit\'e des donn\'ees et de suivre les performances du mod\`ele.

\begin{codeblock}[title=\textbf{G\'en\'erer un rapport de drift avec Evidently}]
\begin{minted}{python}
import pandas as pd
from evidently import ColumnMapping
from evidently.report import Report
from evidently.metric_preset import (
    DataDriftPreset,
    DataQualityPreset,
    TargetDriftPreset,
)


# Charger les donnees
reference_data = pd.read_csv("data/reference.csv")
current_data = pd.read_csv("data/production_week42.csv")

# Definir le mapping des colonnes
column_mapping = ColumnMapping(
    target="label",
    prediction="prediction",
    numerical_features=["age", "income", "credit_score"],
    categorical_features=["region", "product_type"],
)

# Creer le rapport de drift
drift_report = Report(metrics=[
    DataDriftPreset(),
    DataQualityPreset(),
    TargetDriftPreset(),
])

drift_report.run(
    reference_data=reference_data,
    current_data=current_data,
    column_mapping=column_mapping,
)

# Sauvegarder en HTML
drift_report.save_html("reports/drift_report_week42.html")
\end{minted}
\end{codeblock}

\begin{codeblock}[title=\textbf{Tests automatis\'es avec Evidently}]
\begin{minted}{python}
from evidently.test_suite import TestSuite
from evidently.test_preset import (
    DataDriftTestPreset,
    DataStabilityTestPreset,
    NoTargetPerformanceTestPreset,
)
from evidently.tests import (
    TestColumnDrift,
    TestShareOfDriftedColumns,
    TestMeanInNSigmas,
)


# Suite de tests pour la CI
test_suite = TestSuite(tests=[
    # Preset : tester le drift sur toutes les colonnes
    DataDriftTestPreset(),

    # Tests specifiques
    TestColumnDrift(column_name="income", stattest="ks"),
    TestShareOfDriftedColumns(lt=0.3),  # < 30% de drift
    TestMeanInNSigmas(column_name="age", n=2),
])

test_suite.run(
    reference_data=reference_data,
    current_data=current_data,
    column_mapping=column_mapping,
)

# Verifier si les tests passent
result = test_suite.as_dict()
all_passed = all(
    t["status"] == "SUCCESS"
    for t in result["tests"]
)

if not all_passed:
    failed = [
        t["name"] for t in result["tests"]
        if t["status"] == "FAIL"
    ]
    raise RuntimeError(
        f"Tests de drift echoues : {failed}"
    )

print("Tous les tests de monitoring sont passes.")
\end{minted}
\end{codeblock}

%% ----------------------------------------------------------------
\section{Prometheus et Grafana pour la collecte de m\'etriques}
\label{sec:prometheus-grafana}

Prometheus est un syst\`eme de monitoring et d'alerte open-source qui collecte des m\'etriques sous forme de s\'eries temporelles. Grafana est un outil de visualisation qui permet de cr\'eer des tableaux de bord \`a partir de ces m\'etriques.

\begin{codeblock}[title=\textbf{Exposer des m\'etriques ML avec Prometheus}]
\begin{minted}{python}
from prometheus_client import (
    Counter, Histogram, Gauge, Summary,
    start_http_server,
)
import time
import numpy as np

# Definir les metriques
PREDICTION_COUNT = Counter(
    "model_predictions_total",
    "Nombre total de predictions",
    ["model_version", "outcome"],
)

PREDICTION_LATENCY = Histogram(
    "model_prediction_latency_seconds",
    "Latence des predictions en secondes",
    buckets=[0.01, 0.025, 0.05, 0.1, 0.25, 0.5, 1.0],
)

FEATURE_DRIFT_SCORE = Gauge(
    "model_feature_drift_psi",
    "Score PSI de drift par feature",
    ["feature_name"],
)

PREDICTION_CONFIDENCE = Summary(
    "model_prediction_confidence",
    "Distribution de la confiance des predictions",
)


class MonitoredModel:
    """Wrapper qui ajoute le monitoring a un modele."""

    def __init__(self, model, version: str = "v1"):
        self.model = model
        self.version = version

    def predict(self, features: np.ndarray) -> np.ndarray:
        start = time.perf_counter()

        predictions = self.model.predict(features)
        probabilities = self.model.predict_proba(features)

        # Enregistrer la latence
        latency = time.perf_counter() - start
        PREDICTION_LATENCY.observe(latency)

        # Enregistrer les predictions
        for pred in predictions:
            PREDICTION_COUNT.labels(
                model_version=self.version,
                outcome=str(pred),
            ).inc()

        # Enregistrer la confiance
        for prob in probabilities.max(axis=1):
            PREDICTION_CONFIDENCE.observe(prob)

        return predictions

    def update_drift_metrics(
        self, feature_name: str, psi_value: float
    ):
        """Met a jour le score de drift pour une feature."""
        FEATURE_DRIFT_SCORE.labels(
            feature_name=feature_name
        ).set(psi_value)


# Demarrer le serveur de metriques sur le port 8000
start_http_server(8000)
\end{minted}
\end{codeblock}

\begin{codeblock}[title=\textbf{Configuration Prometheus \file{prometheus.yml}}]
\begin{minted}{yaml}
global:
  scrape_interval: 15s
  evaluation_interval: 15s

rule_files:
  - "ml_alerts.yml"

scrape_configs:
  - job_name: "ml-model"
    static_configs:
      - targets: ["ml-model-service:8000"]
    metrics_path: /metrics
    scrape_interval: 10s

  - job_name: "ml-drift-checker"
    static_configs:
      - targets: ["drift-service:8001"]
    scrape_interval: 300s  # Toutes les 5 minutes
\end{minted}
\end{codeblock}

%% ----------------------------------------------------------------
\section{Alertes et tableaux de bord}
\label{sec:alerts}

\begin{codeblock}[title=\textbf{R\`egles d'alerte Prometheus \file{ml\_alerts.yml}}]
\begin{minted}{yaml}
groups:
  - name: ml-model-alerts
    rules:
      - alert: HighPredictionLatency
        expr: |
          histogram_quantile(0.95,
            rate(model_prediction_latency_seconds_bucket[5m])
          ) > 0.5
        for: 10m
        labels:
          severity: warning
        annotations:
          summary: "Latence p95 > 500ms depuis 10min"
          description: >
            La latence au 95e percentile est de
            {{ $value | humanizeDuration }}.

      - alert: DataDriftDetected
        expr: model_feature_drift_psi > 0.2
        for: 30m
        labels:
          severity: critical
        annotations:
          summary: "Drift detecte sur {{ $labels.feature_name }}"
          description: >
            PSI = {{ $value }} (seuil: 0.2).
            Re-entrainement recommande.

      - alert: LowPredictionVolume
        expr: |
          rate(model_predictions_total[1h]) < 10
        for: 30m
        labels:
          severity: warning
        annotations:
          summary: "Volume de predictions anormalement bas"

      - alert: LowModelConfidence
        expr: |
          model_prediction_confidence{quantile="0.5"} < 0.6
        for: 15m
        labels:
          severity: warning
        annotations:
          summary: "Confiance mediane < 60%"
\end{minted}
\end{codeblock}

\begin{bonnepratique}[title=\textbf{Tableau de bord Grafana -- m\'etriques essentielles}]
Un bon tableau de bord ML doit inclure les panneaux suivants~:
\begin{enumerate}
    \item \textbf{Volume de pr\'edictions}~: nombre de requ\^etes par minute (d\'etecter les anomalies de trafic).
    \item \textbf{Distribution des pr\'edictions}~: r\'epartition des classes pr\'edites (d\'etecter le \emph{prior probability shift}).
    \item \textbf{Latence}~: p50, p95, p99 de la latence d'inf\'erence.
    \item \textbf{Confiance}~: distribution des scores de confiance du mod\`ele.
    \item \textbf{Drift par feature}~: score PSI pour chaque variable d'entr\'ee.
    \item \textbf{Performance}~: accuracy, F1-score glissant (si les labels sont disponibles).
    \item \textbf{Erreurs syst\`eme}~: taux d'erreur HTTP, utilisation CPU/m\'emoire.
\end{enumerate}
\end{bonnepratique}

%% ----------------------------------------------------------------
\section{Architecture de monitoring}
\label{sec:monitoring-architecture}

\begin{figure}[htbp]
\centering
\begin{tikzpicture}[
    node distance=1.5cm and 2cm,
    block/.style={rectangle, draw, rounded corners, fill=blue!8,
        minimum width=2.6cm, minimum height=1cm, align=center,
        font=\small},
    store/.style={cylinder, draw, fill=green!8, shape border rotate=90,
        minimum width=1.8cm, minimum height=1cm, aspect=0.3,
        align=center, font=\small},
    alert/.style={rectangle, draw, rounded corners, fill=red!10,
        minimum width=2.2cm, minimum height=0.8cm, align=center,
        font=\small},
    arrow/.style={-{Stealth[length=2.5mm]}, thick},
    darrow/.style={-{Stealth[length=2.5mm]}, thick, dashed}
]

% Application layer
\node[block] (api) {API de\\pr\'ediction};
\node[block, right=of api] (exporter) {Exporter\\m\'etriques};
\node[block, below=of api] (drift) {Service de\\d\'etection drift};
\node[block, left=of api] (users) {Utilisateurs /\\Applications};

% Storage layer
\node[store, right=3cm of drift] (prom) {Prometheus\\(TSDB)};
\node[store, below=of drift] (logs) {Feature\\Store / Logs};

% Visualization and alerts
\node[block, right=of prom] (grafana) {Grafana\\Dashboards};
\node[alert, above=of grafana] (alertmgr) {Alert\\Manager};
\node[alert, right=of alertmgr] (notif) {Slack /\\PagerDuty};

% Re-training
\node[block, below=of logs] (retrain) {Pipeline de\\r\'e-entra\^inement};

% Arrows
\draw[arrow] (users) -- (api);
\draw[arrow] (api) -- (exporter);
\draw[arrow] (exporter) -- (prom);
\draw[arrow] (api) -- (drift);
\draw[arrow] (drift) -- (prom);
\draw[arrow] (api.south) -- ++(0,-0.4) -| (logs);
\draw[arrow] (prom) -- (grafana);
\draw[arrow] (prom) -- (alertmgr);
\draw[arrow] (alertmgr) -- (notif);
\draw[darrow] (drift) -- (logs);
\draw[darrow, red!60] (alertmgr.south) -- ++(0,-0.5) -| (retrain)
    node[pos=0.25, right, font=\scriptsize] {d\'eclencher};
\draw[darrow] (logs) -- (retrain);

\end{tikzpicture}
\caption{Architecture de monitoring pour un mod\`ele ML en production.}
\label{fig:monitoring-architecture}
\end{figure}

%% ----------------------------------------------------------------
\section{Comparaison des outils de monitoring ML}
\label{sec:monitoring-comparison}

\begin{formules}[title=\textbf{Evidently vs WhyLabs vs Fiddler vs NannyML}]
\begin{center}
\begin{tabular}{lcccc}
\toprule
\textbf{Crit\`ere} & \textbf{Evidently} & \textbf{WhyLabs} & \textbf{Fiddler} & \textbf{NannyML} \\
\midrule
Licence & Open-source & Freemium & Commercial & Open-source \\
D\'etection drift & Oui & Oui & Oui & Oui \\
Qualit\'e donn\'ees & Oui & Oui & Limit\'ee & Non \\
Explicabilit\'e & Limit\'ee & Non & Oui & Non \\
Perf. sans labels & Non & Oui & Oui & Oui \\
Int\'egration CI & Oui & Oui & Non & Oui \\
Dashboard & HTML/Cloud & Cloud & Cloud & Limit\'e \\
Complexit\'e & Faible & Mod\'er\'ee & \'Elev\'ee & Faible \\
\bottomrule
\end{tabular}
\end{center}
\end{formules}

\begin{remarque}
NannyML se distingue par sa capacit\'e \`a estimer la performance d'un mod\`ele \emph{sans avoir acc\`es aux labels r\'eels}, gr\^ace \`a la m\'ethode CBPE (\emph{Confidence-Based Performance Estimation}). Cela est particuli\`erement utile lorsque les labels arrivent avec un d\'elai important (ex.~: d\'etection de fraude, pr\'ediction de churn \`a 90 jours).
\end{remarque}

\begin{toolbox}[title=\textbf{Stack de monitoring recommand\'ee}]
Pour un projet MLOps complet, la stack suivante est recommand\'ee~:
\begin{itemize}
    \item \textbf{Evidently} pour la d\'etection de drift et les tests en CI
    \item \textbf{Prometheus} pour la collecte de m\'etriques en temps r\'eel
    \item \textbf{Grafana} pour la visualisation et les tableaux de bord
    \item \textbf{AlertManager} pour les notifications (Slack, PagerDuty, email)
    \item \textbf{NannyML} pour l'estimation de performance sans labels
\end{itemize}
Cette combinaison offre une couverture compl\`ete tout en restant open-source.
\end{toolbox}

%% ----------------------------------------------------------------
\section{Mini-projet~: ajouter le monitoring au mod\`ele d\'eploy\'e}
\label{sec:monitoring-miniproject}

\begin{miniprojet}[title=\textbf{Monitoring du mod\`ele en production}]
\textbf{Objectif}~: ajouter une couche de monitoring compl\`ete au mod\`ele d\'eploy\'e lors des chapitres pr\'ec\'edents.

\textbf{\'Etapes}~:
\begin{enumerate}
    \item \textbf{Instrumenter l'API}~: ajouter les m\'etriques Prometheus (compteur de pr\'edictions, histogramme de latence, jauge de drift) \`a l'API FastAPI/Flask.
    \item \textbf{Logger les pr\'edictions}~: stocker chaque pr\'ediction avec ses features d'entr\'ee dans un fichier CSV ou une base de donn\'ees.
    \item \textbf{Script de drift}~: cr\'eer un script qui compare les donn\'ees de r\'ef\'erence aux donn\'ees de production de la derni\`ere semaine (PSI + KS test).
    \item \textbf{Rapport Evidently}~: g\'en\'erer un rapport HTML de drift hebdomadaire automatis\'e via un cron job ou un workflow GitHub Actions.
    \item \textbf{Prometheus + Grafana}~: d\'eployer Prometheus et Grafana via Docker Compose. Cr\'eer un dashboard avec au moins 5 panneaux.
    \item \textbf{Alertes}~: configurer au moins 3 alertes (latence \'elev\'ee, drift d\'etect\'e, volume anormal).
\end{enumerate}

\textbf{Docker Compose}~:
\begin{minted}{yaml}
version: "3.8"
services:
  ml-model:
    build: .
    ports:
      - "8080:8080"
      - "8000:8000"  # Metriques Prometheus

  prometheus:
    image: prom/prometheus:latest
    volumes:
      - ./monitoring/prometheus.yml:/etc/prometheus/prometheus.yml
      - ./monitoring/ml_alerts.yml:/etc/prometheus/ml_alerts.yml
    ports:
      - "9090:9090"

  grafana:
    image: grafana/grafana:latest
    ports:
      - "3000:3000"
    environment:
      - GF_SECURITY_ADMIN_PASSWORD=admin
    volumes:
      - grafana-data:/var/lib/grafana

volumes:
  grafana-data:
\end{minted}

\textbf{Livrables}~:
\begin{itemize}
    \item Code source de l'API instrument\'ee avec les m\'etriques Prometheus.
    \item Fichier \file{docker-compose.yml} pour le d\'eploiement complet.
    \item Capture d'\'ecran du dashboard Grafana avec les m\'etriques ML.
    \item Rapport Evidently g\'en\'er\'e automatiquement.
    \item Documentation des alertes configur\'ees.
\end{itemize}
\end{miniprojet}

%% ----------------------------------------------------------------
\section{Exercices}
\label{sec:monitoring-exercises}

\begin{exercice}[Niveau 1 -- D\'etection de drift basique]
\begin{enumerate}
    \item G\'en\'erez deux distributions normales~: $X_{\text{ref}} \sim \mathcal{N}(0, 1)$ avec 10\,000 \'echantillons et $X_{\text{cur}} \sim \mathcal{N}(0.5, 1.2)$ avec 10\,000 \'echantillons.
    \item Calculez le PSI entre les deux distributions. Interpr\'etez le r\'esultat.
    \item Appliquez le test de Kolmogorov-Smirnov. Le drift est-il d\'etect\'e avec $\alpha = 0.05$~?
    \item Calculez la divergence KL dans les deux sens. Observez l'asym\'etrie.
    \item Utilisez Evidently pour g\'en\'erer un rapport de drift sur ces donn\'ees synth\'etiques.
\end{enumerate}
\end{exercice}

\begin{exercice}[Niveau 2 -- Monitoring d'une API de pr\'ediction]
\begin{enumerate}
    \item Cr\'eez une API FastAPI qui sert un mod\`ele scikit-learn (classificateur de votre choix).
    \item Ajoutez les m\'etriques Prometheus suivantes~: nombre de pr\'edictions, latence (histogramme), distribution des classes pr\'edites.
    \item D\'eployez Prometheus et Grafana avec Docker Compose.
    \item Cr\'eez un dashboard Grafana affichant les m\'etriques en temps r\'eel.
    \item \'Ecrivez un script de charge qui envoie des requ\^etes avec des donn\'ees progressivement drift\'ees.
    \item Configurez une alerte qui se d\'eclenche lorsque la latence p95 d\'epasse 200\,ms.
\end{enumerate}
\end{exercice}

\begin{exercice}[Niveau 3 -- Syst\`eme de monitoring autonome avec re-entra\^inement]
Concevez et impl\'ementez un syst\`eme de monitoring complet~:
\begin{enumerate}
    \item \textbf{D\'etection multi-m\'ethode}~: impl\'ementez PSI, KS, et la m\'ethode ADWIN (\emph{Adaptive Windowing}) pour la d\'etection de drift en temps r\'eel.
    \item \textbf{Pipeline de monitoring}~: cr\'eez un service qui~:
    \begin{itemize}
        \item Collecte les pr\'edictions et features toutes les 5 minutes
        \item Calcule les scores de drift sur une fen\^etre glissante
        \item Met \`a jour les m\'etriques Prometheus
    \end{itemize}
    \item \textbf{Re-entra\^inement automatique}~: lorsqu'un drift significatif est d\'etect\'e (PSI $> 0.2$ pendant plus de 30 minutes), d\'eclenchez automatiquement un pipeline de re-entra\^inement via l'API GitHub Actions.
    \item \textbf{Validation post-entra\^inement}~: le nouveau mod\`ele doit passer les tests de performance avant d\'eploiement (\emph{champion/challenger}).
    \item \textbf{Observabilit\'e compl\`ete}~: dashboard Grafana avec au moins 10 panneaux couvrant performance, drift, infrastructure et m\'etriques m\'etier.
\end{enumerate}
\end{exercice}

%% ----------------------------------------------------------------
\section{Aide-m\'emoire}
\label{sec:monitoring-cheatsheet}

\begin{cheatsheet}[title=\textbf{Monitoring des mod\`eles ML}]
\textbf{Types de drift}~:
\begin{itemize}
    \item \textbf{Covariate shift}~: $P(X)$ change, $P(Y|X)$ stable
    \item \textbf{Prior probability shift}~: $P(Y)$ change
    \item \textbf{Concept drift}~: $P(Y|X)$ change (le plus grave)
\end{itemize}

\textbf{M\'etriques de drift}~:
\begin{itemize}
    \item \textbf{PSI}~: $< 0.1$ OK, $0.1$--$0.2$ surveiller, $\geq 0.2$ agir
    \item \textbf{KS test}~: $p < 0.05 \Rightarrow$ drift d\'etect\'e
    \item \textbf{KL divergence}~: asym\'etrique, $\geq 0$, $= 0$ si identiques
\end{itemize}

\textbf{Code Evidently essentiel}~:
\begin{minted}{python}
from evidently.report import Report
from evidently.metric_preset import DataDriftPreset

report = Report(metrics=[DataDriftPreset()])
report.run(
    reference_data=ref_df,
    current_data=cur_df,
)
report.save_html("drift_report.html")
\end{minted}

\textbf{Prometheus -- m\'etriques cl\'es}~:
\begin{minted}{python}
from prometheus_client import Counter, Histogram, Gauge
PREDICTIONS = Counter("predictions_total", "...", ["class"])
LATENCY = Histogram("prediction_latency_s", "...")
DRIFT = Gauge("feature_drift_psi", "...", ["feature"])
\end{minted}

\textbf{Checklist monitoring}~:
\begin{enumerate}
    \item[$\square$] M\'etriques Prometheus expos\'ees par l'API
    \item[$\square$] Pr\'edictions et features logu\'ees
    \item[$\square$] D\'etection de drift automatis\'ee
    \item[$\square$] Dashboard Grafana op\'erationnel
    \item[$\square$] Alertes configur\'ees (latence, drift, volume)
    \item[$\square$] Processus de re-entra\^inement d\'efini
    \item[$\square$] Strat\'egie de rollback pr\^ete
\end{enumerate}
\end{cheatsheet}
